\section{Stetigkeit und Funktionen}


\Title{Definitionen}

\begin{tcolorbox}
    \Def Sei $D \subseteq \R,$ $x_0 \in D$. Die Funktion $f : D \to \R$ ist \textbf{stetig 
    in einem Punkt} $x_0$, falls es für jedes $\epsilon > 0$ ein $\delta > 0$ gibt, so dass 
    für alle $x \in D$ gilt:
    $$|x - x_0| < \delta \Rightarrow |f(x) - f(x_0)| < \epsilon$$

    \textbf{Alternativ:} Falls für jede Folge $(a_n)_{n \geq 1}$ mit $\Limit
    a_n = x_0$ folgendes gilt:
    $$f(\Limit a_n) = f(x_0) = \Limit f(a_n)$$
    ist die Funktion $f$ in $x_0$ stetig.
\end{tcolorbox}

\Def Die Funktion $f : D \to \R$ ist stetig, falls sie in jedem Punkt von $D$ stetig ist.

\vspace{2mm}
\Bsp Abrundungsfunktion $\left\lceil \cdot \right\rceil$  ist in jedem Punkt $x_{0} \notin \mathbb{Z}$ 
stetig, aber in keinem Punkt $x \in \mathbb{Z}$.

\begin{tcolorbox}
    \textbf{Gleichmässige Stetigkeit}: Eine Funktion $f : D \to \R$ ist gleichmässig stetig falls:
    \begin{center}
        $\forall \epsilon > 0$, $\exists \delta > 0$, $\forall x,y \in D$: $|x - y| < \delta \Rightarrow |f(x) - f(y)| < \epsilon$
    \end{center}
    Wobei Funktionen welche in einem kompakten Intervall stetig sind, im selben Intervall auch gleichmässig
    stetig sind.
\end{tcolorbox}


\Title{Rechnen mit Stetigkeit}

Sei $x_0 \in D \subseteq \R$, $\lambda \in \R$ und $f : D \to \R$, $g : D \to \R$ beide in $x_0$ stetig.
So gilt:

\begin{enumerate}[label=(\arabic*)]
    \item $f + g$, $\lambda \cdot f$, $f \cdot g$ sind stetig in $x_0$ (auch $g \circ f$ bei anderen Def. Bereichen)
    \item falls $g(x_0) \neq 0$, dann ist $\frac{f}{g}$ stetig in $x_0$ für $D \cap \{ x \in D : g(x) \neq 0 \}$
    \item $|f|$, $\max(f,g)$ und $\min(f,g)$ sind stetig in $x_0$
    \item Polynomiale Funktionen sind auf ganz $\R$ stetig
    \item die Trigonometrischen Funktionen $\sin : \R \to \R$ und $\cos : \R \to \R$ sind stetig
    \item die Exponentialfunktion $e^x$ ist auf ganz $\R$ stetig
    \item seien $P, Q$ polynomiale Funktionen auf $\R$ mit $Q \neq 0$. Seien $x_1, ... , x_m$ die
    Nullstellen von $Q$. Dann ist $\frac{P}{Q}$ stetig für $\R \backslash \{x_1, ... , x_m\}$ 
    (komplexe Nullstellen sind immer komplementär)
    \item sei $f : D_1 \to D_2$ und $g : D_2 \to \R$ Funktionen, sowie $x_0 \in D_1$. Falls $f$ in $x_0$
    und $g$ in $f(x_0)$ stetig sind, so ist $g(f(x_0)) : D_1 \to \R$ in $x_0$ stetig
\end{enumerate}


\Title{Zwischenwertsatz}

\begin{tcolorbox}
    \Satz Sei $I \subseteq \R$ ein Intervall, $f : I \to \R$ eine stetige Funktion und $a,b \in I$.
    Für jedes $c$ zwischen $f(a)$ und $f(b)$ gibt es ein $z$ zwischen $a$ und $b$ mit $f(z) = c$.
\end{tcolorbox}

Der Zwischenwertsatz wird oftmals verwendet um zu zeigen, dass eine Funktion einen gewissen Wert annimmt.
Typische Anwendungsbeispiele sind:
\begin{enumerate}[label=(\arabic*)]
    \item Sei $f : [a,b] \to \R$ stetig. Falls $f(a) \cdot f(b) < 0$, dann $\exists c \in ]a,b[$ mit
    $f(c) = 0$ (also eine Nullstelle).
    \item Sei $P(x) = a_nx^n + ... + a_0$ ein Polynom mit $a_n \neq 0$ und $n$ ungerade. Dann besitzt
    $P$ mindestens eine Nullstelle in $R$.
\end{enumerate}


\Title{Min-Max Satz}

\Def Ein Intervall $I \subseteq \R$ ist \textbf{kompakt}, falls es von der Form \text{[a,b]}  
ist mit $a \leq b$.

\Bem Für $x_0$ in einem kompakten Intervall gibt es immer min. eine Folge $(a_n)_{n \geq 1}$, $a_n \in I$, 
so dass $\Limit a_n = x_0$.

\begin{tcolorbox}
    \Satz Sei $f : I = [a,b] \to \R$ stetig auf einem kompakten Intervall $I$. Dann gibt es 
    $u, v \in I$ mit
    $$f(u) \leq f(x) \leq f(v) \quad \quad \forall x \in I.$$
    Insbesondere ist $f$ beschränkt.
\end{tcolorbox}


\columnbreak
%----------------------------------------------------------------------------------------------------	


\Title{Satz der Umkehrabbildung}

\Satz Sei $I \subseteq \R$ ein Intervall und $f : I \to \R$ eine stetige, streng monotone Funktion.
Dann ist das Bild $f(I) := J$ ein Intervall und die Umkehrfunktion $f^{-1} : J \to I$ ist stetig,
streng monoton.

\Bsp Sei $n \geq 1$. Dann ist $f : [0, \infty[ \to [0, \infty[$ als $x \to x^n$ streng monoton wachsend,
stetig und surjektiv. Nach dem Umkehrsatz existiert nun  eine streng monoton wachsende, stetige Umkehrabbildung
$f^{-1} : [0, \infty[ \to [0, \infty[$ als $x \to \sqrt[n]{x}$.


\Title{Konvergenz von Funktionenfolgen}

\Def Eine \textbf{Funktionenfolge} ist eine Abbildung $\N \to \R^D : n \to f_n$, wobei $f_n$ die n-te Funktion
ist. Für jedes $x \in D$ erhält man eine Folge $(f_n(x))_{n \geq 0}$ in $\R$.

\begin{tcolorbox}
    Die Funktionenfolge  $(f_n)_{n \geq 0}$ \textbf{konvergiert punktweise} gegen eine Funktion 
    $f : D \to \R$, falls für alle $x \in D$:
    $$f(x) = \Limit{f_n(x)}$$
    \textbf{Alternativ:}
    $$\forall \epsilon > 0, \; \forall x \in D, \; \exists N \geq 1, \; \forall n \geq N, : |f_n(x) - f(x)| < \epsilon$$
\end{tcolorbox}


\begin{tcolorbox}
    Die Funktionenfolge \textbf{konvergiert gleichmässig} gegen $f$ falls:
    $$\forall \epsilon > 0, \; \exists N \geq 1, \; \forall n \geq N, \; \forall x \in D : |f_n(x) - f(x)| < \epsilon$$
    \textbf{Alternativ:}
    $$\forall \epsilon > 0, \; \exists N \geq 1, \; \forall n,m \geq N, \; \forall x \in D : |f_n(x) - f_m(x)| < \epsilon$$
\end{tcolorbox}

\Bem Jede gleichmässig konvergente Funktionenfolge ist auch punktweise konvergent, jedoch nicht 
umgekehrt!

\Satz Sei $D \subseteq \R$ und $f_n : D \to \R$ eine Funktionenfolge aus stetigen Funktionen.
Wenn $f_n$ gleichmässig gegen $f$ konvergiert, dann ist auch $f$ stetig. (Für punktweise konvergenz
muss dies nicht der Fall sein.)


\Title{Strategie - Konvergenz von Funktionenfolgen}

\begin{enumerate}[label=(\arabic*)]
    \item Punktweiser Limes von $f_n$ auf $D$ finden.
    \item Prüfe auf gleichmässige Konvergenz:
        \begin{enumerate}[label=(\roman*)]
            \item Indirekte Methode: $f$ unstetig bedeuted keine gleichmässige Konvergenz, $f$ stetig,
            monoton wachsend und $D$ kompakt bedeuted gleichmässige Konvergenz.
            \item Direkte Methode: Berechne $\sup_{x \in D} |f_n(x) - f(x)|$ (evtl. Ableitung von 
            $|f_n(x) - f(x)|$ gleich Null setzen). Danach Limes für $n \to \infty$ von $\sup_{x \in D} |f_n(x) - f(x)|$
            berechnen, falls dieser Null ist, so konvergiert $f_n$ gleichmässig.
        \end{enumerate}
\end{enumerate}


\Title{Konvergenz von Funktionenreihen}

\Def Die Reihe $\Reihe f_k(x)$ \textbf{konvergiert gleichmässig}, falls die durch $S_n(x) = \sum_{k=0}^{n} f_k(x)$
definierte Funktionenfolge gleichmässig konvergiert.

\begin{tcolorbox}
    \Satz Sei $f_n$ eine Folge stetiger Funktionen und es gilt $|f_n(x)| \leq c_k, \forall x \in D$.
    Wenn nun $\Reihe c_k$ konvergiert, so konvergiert auch die Reihe $\Reihe f_k(x)$ gleichmässig,
    wobei der Grenzwert $f(x)$ eine stetige Funktion ist.
\end{tcolorbox}


\columnbreak
%----------------------------------------------------------------------------------------------------


\Title{Grenzwerte von Funktionen}

Wir betrachten Funktionen $f : D \to \R$ und wollen den Grenzwert für den Fall $x \in D$ strebt gegen
$x_0$ finden. Dabei kann es sein, dass $x_0 \notin D$. Wir nehmen an, dass $x_0$ ein Häufungspunkt 
von $D$ ist.

\begin{tcolorbox}
    $x_0 \in \R$ ist ein \textbf{Häufungspunkt} der Menge $D$ falls:
    $$\forall \delta > 0 : \quad (]x_0 - \delta, x_0 + \delta [ \backslash \{ x_0 \}) \cap D \neq \emptyset$$
\end{tcolorbox}

Nun können wir den Grenzwert von Funktionen betrachten.

\begin{tcolorbox}
    Sei $f$ eine Funktion und $x_0$ ein Häufungspunkt von $D$. Dann ist $A \in \R$ der Grenzwert
    von $f(x)$ gegen $x_0$, falls $\forall \epsilon > 0, \exists \delta > 0$ so dass
    $$\forall x \in D \cap (]x_0 - \delta, x_0 + \delta [ \backslash \{ x_0 \}) : |f(x) - A| < \epsilon$$
\end{tcolorbox}

\Satz Sei $(a_n)_{n \geq 1}, \, a_n \in \C$ beschränkt:
\begin{enumerate}[label=(\arabic*)]
    \item{Es gibt mindestens einen Häufungspunkt der Folge}
    \item{Falls es nur einen Häufungspunkt \(a\in\C\) gibt, konvergiert \((a_n)_{n \geq 1}\) gegen \(a\).}
\end{enumerate}

\Bem Seien $f,g$ Funktionen und $x_0$ ein Häufungspunkt von $D$ so gilt:
\begin{enumerate}[label = (\arabic*)]
    \item $\lim_{x \to x_0} f(x) = A$, wenn für jede Folge $(a_n)_{n \geq 1}$ in $D \backslash \{ x_0 \}$
    mit $\lim_{n \to \infty} a_n = x_0$ folgt $\lim_{n \to \infty} f(a_n) = A$
    \item  Wenn $x_0 \in D$, dann ist $f$ stetig in $x_0$ falls $\lim_{x \to x_0} f(x) = f(x_0)$
    \item $\lim_{x \to x_0} (f + g)(x) = \lim_{x \to x_0} f(x) + \lim_{x \to x_0} g(x)$
    \item $\lim_{x \to x_0} (f \cdot g)(x) = \lim_{x \to x_0} f(x) \cdot \lim_{x \to x_0} g(x)$
    \item Sei $f \leq g$, so ist $\lim_{x \to x_0} f(x) \leq \lim_{x \to x_0} g(x)$
    \item Falls $g_1 \leq f \leq g_2$ und $\lim_{x \to x_0} g_1(x) = \lim_{x \to x_0} g_2(x)$, so
    existiert $\lim_{x \to x_0} f(x) = \lim_{x \to x_0} g_1(x)$
\end{enumerate}

\Satz Sei $x_0$ ein Häufungspunkt von $D$, $f : D \to E$ eine Funktion mit $y_0 = \Lim_{x \to x_0} f(x)$
mit $y_0$. Fall $g : E \to \R$ stetig ist in $y_0$ folgt:
$$\Lim_{x \to x_0} g(f(x)) = g(y_0)$$


\Title{Linksseitiger und Rechtsseitiger Grenzwert}

Sei $f : D \to \R$ und $x_0 \in D$ ist ein Häufungspunkt von $D \; \cap \; ]x_0, + \infty[$. Falls der
Grenzwert der eingeschränkten Funktion $f$ im Bereich $D \; \cap \; ]x_0, + \infty[$ für $x \to x_0$
existiert, wird er mit $\lim_{x \to x_0^+} f(x)$ bezeichnet und nennt sich \textbf{rechtsseitiger
Grenzwert} von $f$ bei $x_0$. Das Analoge gilt für den \textbf{linksseitigen Grenzwert}.

Wir erweitern diese Definition auf $\lim_{x \to x_0^+} f(x) = + \infty$ falls gilt:
$$\forall \epsilon > 0, \exists \delta > 0, \forall x \in D \cap ]x_0, x_0 + \delta[ : \; f(x) > \frac{1}{\epsilon}$$
und analog für $\lim_{x \to x_0^+} f(x) = - \infty$:
$$\forall \epsilon > 0, \exists \delta > 0, \forall x \in D \cap ]x_0, x_0 + \delta[ : \; f(x) < -\frac{1}{\epsilon}$$
Für den linksseitigen Grenzwert gilt das Analoge.


\columnbreak
%----------------------------------------------------------------------------------------------------